\exercisesection{分类}

\begin{enumerate}
\def\labelenumi{\arabic{enumi})}
\tightlist
\item
  \begin{enumerate}
  \def\labelenumii{\alph{enumii})}
  \tightlist
  \item
    为使 G 具有可数基，必要且充分的是 \(\widehat{G}\) 具有可数基。（若 G 有可数基，则 \(L^{1}(G)\) 有可数基，从而 \(\mathsf{X}(\mathsf{L}^{1}(\mathsf{G}))\) 有可数基。）
  \end{enumerate}
\end{enumerate}

\begin{enumerate}
\def\labelenumi{\alph{enumi})}
\setcounter{enumi}{1}
\item
  为使 G 在无穷远处可数，必要且充分的是 \(\widehat{\mathrm{G}}\) 容有一个具有可数基的开子群，或者等价地，\(\widehat{\mathrm{G}}\) 是可度量化的。（利用 II，第 248 页，命题 3。）
\item
  设 G 是紧的。为使 G 可度量化，必要且充分的是 \(\widehat{G}\) 是可数的，或者等价地，G 同构于 \(T^{N}\) 的一个闭子群。（一个可数的离散交换群是 \(\mathbf{Z}^{(N)}\) 的一个商。）
\end{enumerate}

\begin{enumerate}
\def\labelenumi{\arabic{enumi})}
\setcounter{enumi}{1}
\tightlist
\item
  \begin{enumerate}
  \def\labelenumii{\alph{enumii})}
  \tightlist
  \item
    每个无限交换群 \(\Gamma\) 都容许一个无限的可数商群。（利用 A，II，第 185 页，练习 14，把一个无限可数子群 \(\Delta\)（\(\Gamma\) 的）嵌入到一个可数的可除群 \(\Delta'\) 中；然后在 \(\Gamma\) 上延拓 \(\Delta\) 到 \(\Delta'\) 的恒等同态。）
  \end{enumerate}
\end{enumerate}

\begin{enumerate}
\def\labelenumi{\alph{enumi})}
\setcounter{enumi}{1}
\tightlist
\item
  每个无限的紧交换群都包含一个无限的可度量化闭子群。（利用 a)。）
\end{enumerate}

\begin{enumerate}
\def\labelenumi{\arabic{enumi})}
\setcounter{enumi}{2}
\tightlist
\item
  \begin{enumerate}
  \def\labelenumii{\alph{enumii})}
  \tightlist
  \item
    设 K 为紧群，\(\mathrm{K}_{(n)}\) 为 K 中阶为 n 的元素所成的集合。若对每个 n 都有 \(K \neq E_{n}\)，则 K 中无限阶元素所成的集合在 K 中稠密。（若 \(\mathrm{K}_{(n)}\) 有非空内部，则有 \(\mathrm{K}_{(m)} = \mathrm{K}\)（对某个整数 \(m \geqslant n\)）。）
  \end{enumerate}
\end{enumerate}

\begin{enumerate}
\def\labelenumi{\alph{enumi})}
\setcounter{enumi}{1}
\item
  设 \(\Gamma\) 为一个离散交换群，其元素的阶不是有界的。则 \(\Gamma\) 容许一个具有同样性质的可数商群。（按练习 2 的 a) 同样论证。）
\item
  若 e 在 G 中的每个邻域都包含一个无限阶元素，则 G 具有一个带同样性质的可度量化闭子群。（化归到 G 是紧的情形，并利用 a) 与 b)。）反之，若 G 不是离散的，则存在整数 q \textgreater{} 1，使得 G 包含一个同构于 \((\mathbf{Z}/q\mathbf{Z})^{\mathbf{N}}\) 的闭子群。（利用 A，VII，第 54 页，练习 4，c)。）
\end{enumerate}

\begin{enumerate}
\def\labelenumi{\arabic{enumi})}
\setcounter{enumi}{3}
\item
  若 G 是紧的且 \(\widehat{G}\) 有一个无限阶元素，则存在一个从 R 到 G 的非平凡连续同态。（由于 R 是可除的，存在一个从 \(\widehat{G}\) 到 R 的非平凡同态。）
\item
  设 p 为素数。群 \(Q_{p}\) 不是紧群与离散群的乘积。（\(Q_{p}\) 的每个紧子群都具有形式 \(p^{n}Z_{p}\)（对某个整数 \(n \in Z\)）。）
\item
  若 G 是挠群，则 G 与 \(\widehat{G}\) 都是全不连通的。（由 II，第 250 页，推论 2，\(\widehat{G}\) 是全不连通的；为证 G 全不连通，利用 II，第 248 页，命题 3 化归到 G 为紧的情形，然后证明 \(\widehat{G}\) 是挠群。）
\item
  我们称元素 \(x\)（属于 \(G\)）为无限高的，如果对每个 \(n \in \mathbf{Z}\)，存在 \(y \in G\) 使得 \(y^n = x\)。证明：若 \(G\) 是紧的，则无限高元素所成的集合是 \(G\) 的中性分支。
\end{enumerate}

¶ 8) 群 G 是局部连通的，当且仅当 G 同构于一个形如 \(\mathbf{R}^{n}\times\mathrm{E}\times\widehat{\mathrm{D}}\) 的群，其中 E 与 D 是离散的，且 D 的每个有限秩子群都是自由的。

\begin{enumerate}
\def\labelenumi{\arabic{enumi})}
\setcounter{enumi}{8}
\tightlist
\item
  \begin{enumerate}
  \def\labelenumii{\alph{enumii})}
  \tightlist
  \item
    设 \(\boldsymbol{a}=(a_{n})_{n\geqslant0}\) 是一个整数序列，均 \textgreater{} 1。我们给 Z 赋予如下的拓扑环结构：其中子集 \(Za_{0}a_{1}\cdots a_{n}\)（\(n\geqslant0\)）构成 0 的一个基本邻域系。我们用 \(\Delta_{a}\) 表示这个环的完备化。它是紧的、可度量的、全不连通的。若 p 是素数且对所有 i 有 \(a_{i}=p\)，则我们有 \(\Delta_{a}=Z_{p}\)。
  \end{enumerate}
\end{enumerate}

\begin{enumerate}
\def\labelenumi{\alph{enumi})}
\setcounter{enumi}{1}
\item
  我们用 \(\mathbf{Z}(\boldsymbol{a}^{\infty})\) 表示形如 \(\exp(2i\pi l/(a_{0}a_{1}\cdots a_{r}))\)（其中 \(l \in Z\) 且 \(r \in N\)）的复数所成的乘法群，并赋予离散拓扑。则 \(\chi \mapsto \chi(1)\) 是从 \(\widehat{\Delta}_{\boldsymbol{a}}\) 到 \(\mathbf{Z}(\boldsymbol{a}^{\infty})\) 上的同构。
\item
  设 B 为 \(\mathbf{R} \times \Delta_{\boldsymbol{a}}\) 的由 \((n, n)\)（\(n \in \mathbf{Z}\)）组成的子群；它是离散的。我们置 \(\Sigma_{\boldsymbol{a}} = (\mathbf{R} \times \Delta_{\boldsymbol{a}})/\mathrm{B}\)。它是一个紧的、可度量的连通群。（注意 \(\mathbf{R}\) 在 \(\Sigma_{\boldsymbol{a}}\) 中的典范像是稠密的，并利用 II，第 244 页，引理 1。）
\end{enumerate}

\(d)\) 设 \(\Gamma_{\boldsymbol{a}}\) 为离散群 \(\mathbf{Q}\) 的由形如 \(l/(a_{0}a_{1}\cdots a_{r})\)（其中 \(l\in\mathbf{Z}\) 且 \(r\in\mathbf{N}\)）的有理数组成的加法子群。若 \(\chi\in\widehat{\Sigma}_{\boldsymbol{a}}\)，则 \(\chi\) 定义了 \(\mathbf{R}\times\Delta_{\boldsymbol{a}}\) 的一个酉特征标，从而定义了 \(\mathbf{R}\) 的一个形如 \(x\mapsto\exp(2i\pi\alpha_{\chi}x)\) 的酉特征标（其中 \(\alpha_{\chi}\in\mathbf{R}\)）。则 \(\chi\mapsto\alpha_{\chi}\) 是从 \(\widehat{\Sigma}_{\boldsymbol{a}}\) 到 \(\Gamma_{\boldsymbol{a}}\) 上的同构。

\begin{enumerate}
\def\labelenumi{\alph{enumi})}
\setcounter{enumi}{4}
\item
  证明 \(\widehat{\mathbf{Q}}\) 典范地同构于 \(\Sigma_{\pmb{a}}\)，其中序列 \(\pmb{a}\) 由 \(a_{n} = n + 2\)（对所有 \(n \geqslant 0\)）定义。
\item
  设 \(R_{d}\) 为赋予离散拓扑的群 R。证明 \(\widehat{R}_{d}\) 同构于 \((\widehat{\Sigma}_{\boldsymbol{a}})^{\mathfrak{c}}\)，其中 \(a_{n}=n+2\)，c 是连续统的基数。（考察 R 作为 Q-向量空间的一组基。）
\end{enumerate}

¶ 10) 称一个拓扑群为单生的，如果存在群中的一个元素，它的幂在群中稠密；称它为螺线管群，如果存在从 R 到该群中的连续同态，其像是稠密的。

\begin{enumerate}
\def\labelenumi{\alph{enumi})}
\item
  沿用练习 9 的记号，群 \(\Delta_{\pmb{a}}\) 是单生的。
\item
  每个紧的单生全不连通群都拓扑同构于形如 \(\Delta_{a}\) 的群。
\item
  设 G 是紧的。为使 G 是单生的，必要且充分的是 \(\widehat{G}\) 同构于赋予离散拓扑的 T 的一个子群。
\item
  设 G 是紧的。下列条件等价：
\end{enumerate}

\begin{enumerate}
\def\labelenumi{(\roman{enumi})}
\item
  G 是螺线管群；
\item
  \(\widehat{\mathbf{G}}\) 同构于赋予离散拓扑的 \(\mathbf{R}\) 的一个子群；
\item
  \(\widehat{\mathbf{G}}\) 是无挠的，且 \(\operatorname{Card}(\widehat{\mathbf{G}}) \leqslant \mathfrak{c}\)（连续统的基数）；
\item
  G 是 \((\Sigma_{\boldsymbol{a}})^{\mathfrak{c}}\) 的一个商，其中 \(\boldsymbol{a}=(n+2)_{n\geqslant0}\)。
\end{enumerate}

\begin{enumerate}
\def\labelenumi{\alph{enumi})}
\setcounter{enumi}{4}
\tightlist
\item
  若 G 是紧的螺线管群，则 G 是单生的。
\end{enumerate}

\begin{enumerate}
\def\labelenumi{\arabic{enumi})}
\setcounter{enumi}{10}
\tightlist
\item
  我们用 L(G) 表示从 G 到 R 中的连续表示所成的集合。G 的单参数子群是指从 R 到 G 中的连续同态的像。
\end{enumerate}

\begin{enumerate}
\def\labelenumi{\alph{enumi})}
\tightlist
\item
  下列条件等价：
\end{enumerate}

\begin{enumerate}
\def\labelenumi{(\roman{enumi})}
\item
  G 是诸单参数子群的并；
\item
  从 Z 到 G 的每个同态都可延拓为从 R 到 G 的连续同态；
\item
  从 \(\widehat{\mathbf{G}}\) 到 \(\mathbf{R} / \mathbf{Z}\) 的每个连续同态都由 \(\mathrm{L}(\widehat{\mathrm{G}})\) 中的一个元素过渡到商而得到；
\item
  \(\widehat{\mathbf{G}}\) 的每个酉特征标都形如 \(e^{i\lambda}\)（其中 \(\lambda \in \mathrm{L}(\widehat{\mathbf{G}})\)）。
\end{enumerate}

\begin{enumerate}
\def\labelenumi{\alph{enumi})}
\setcounter{enumi}{1}
\item
  若 G 有可数基，则 a) 中的条件等价于：G 同构于 \(R^{n} \times T^{I}\)（其中 I 是一个可数集）。
\item
  G 的诸单参数子群的并是 G 的中性分支的一个稠密子群。（利用练习 4。）
\end{enumerate}

\(d)\) 设 \(t \mapsto \chi_t\) 是从 \([0,1]\) 到 \(\widehat{\mathbf{G}}\) 中的连续映射，使得 \(\chi_0\) 是平凡特征标。存在一个映射 \(t \mapsto \lambda_t\)（从 \([0,1]\) 到 \(\mathrm{L}(\mathrm{G})\)，满足 \(\lambda_0 = 0\)），使得对所有 \(t\) 有 \(\chi_t = \exp(2i\pi\lambda_t)\)，并且对每个 \(x \in \mathrm{G}\)，映射 \(t \mapsto \lambda_t(x)\) 是连续的。

\begin{enumerate}
\def\labelenumi{\alph{enumi})}
\setcounter{enumi}{4}
\tightlist
\item
  G 的诸单参数子群的并也是 G 的包含 e 的弧的并，G 的弧是指从 \([0,1]\) 到 G 中的连续映射的像。（利用 d)。
\end{enumerate}

\begin{enumerate}
\def\labelenumi{\arabic{enumi})}
\setcounter{enumi}{11}
\tightlist
\item
  我们沿用练习 11 中的记号 L(G)。
\end{enumerate}

\begin{enumerate}
\def\labelenumi{\alph{enumi})}
\item
  设 H 为 G 的闭子群。L(H) 的每个元素都可延拓为 L(G) 的元素。（先化归到 \(G = R^{n} \times D\)（D 离散）的情形，再化归到 L(H) 的给定元素在 \(R^{n}\) 上为平凡的情形，然后化归到 G 为离散的情形。最后利用 R 是可除的这一事实。）
\item
  从 H 到 \(R^{n} \times T^{I}\)（其中 I 是任意集合）的每个连续态射都可延拓为从 G 到 \(R^{n} \times T^{I}\) 的连续态射。（利用 a) 以及 II，第 226 页，定理 4。）
\item
  反之，设 A 为一个局部紧交换群。假设对每个局部紧交换群 G、对 G 的每个闭子群 H 以及从 H 到 A 的每个连续态射 \(\varphi\)，\(\varphi\) 都可延拓为从 G 到 A 的连续态射。则存在整数 n 与集合 I，使得 A 同构于 \(R^{n} \times T^{I}\)。（取 G = R，H = Z，可以看出 A 是连通的，从而形如 \(R^{n} \times \widehat{D}\)（D 离散）。证明 D 是一个投射 Z-模，从而是自由的。）
\end{enumerate}

\begin{enumerate}
\def\labelenumi{\arabic{enumi})}
\setcounter{enumi}{12}
\tightlist
\item
  设 C 为 G 的紧子集，\(\varepsilon > 0\)。
\end{enumerate}

\begin{enumerate}
\def\labelenumi{\alph{enumi})}
\item
  设 \(\alpha\) 为 G 的 Haar 测度。存在 G 的紧子集 D，使得 \(\alpha(\mathrm{CD}) \leqslant (1 + \varepsilon)\alpha(\mathrm{D})\)。（\(\mathbf{G} = \mathbf{R}^n\) 的情形是直接的，我们化归到 G 容有一个紧开子群 H 的情形，并且可以假设 C 关于 H 是饱和的；于是化归到 G 为离散的情形；接着可以假设 G 是有限生成的，最后 \(\mathbf{G} = \mathbf{Z}^p\)。）
\item
  存在 \(k \in \mathrm{L}^{1}(\widehat{\mathrm{G}})\) 使得 \(\mathcal{F}(k) \in \mathcal{K}(\mathrm{G})\) 且 \(\mathcal{F}(k)(x) = 1\)（对 \(x \in C\)），并有 \(\|k\|_{1} \leqslant 1 + \varepsilon\)。（取形如 \(\lambda fg\) 的 k，其中 \(\lambda\) 是常数，\(f, g \in \mathrm{L}^{2}(\widehat{\mathrm{G}})\) 的 Fourier 变换是适当集合的特征函数，利用 a)。）
\end{enumerate}

¶ 14) a) 假设 e 在 G 中的每个邻域都包含一个无限阶元素。存在 G 的一个紧的、可度量的、完满的全不连通子集 P，它是一个 Kronecker 集（II，第 265 页，练习 15），并且存在一个集中在 P 上的非零弥散测度。（利用练习 3 的 c) 化归到 G 可度量化的情形。然后仿照 TG，IX，第 64 页，引理 3 的构造，并同时利用 II，第 265 页，练习 15 的 e)。）

\begin{enumerate}
\def\labelenumi{\alph{enumi})}
\setcounter{enumi}{1}
\item
  若 \(\mathrm{G} = (\mathbf{Z}/q\mathbf{Z})^{\mathbf{N}}\)，则存在 G 的一个紧的完满全不连通子集 P，它是 \(K_{q}\) 型的，并且存在一个集中在 P 上的非零弥散测度。
\item
  设 G 不是离散的。存在元素 \(\tau\)（属于 \(\mathcal{M}^{1}(G)\)）使得 \(\|\mathcal{F}(\tau)\|_{\infty}<\varrho(\tau)\)，其中 \(\varrho(\tau)\) 是 \(\mathcal{M}^{1}(G)\) 中的谱半径。存在不可逆元素 \(\tau'\)（属于 \(\mathcal{M}^{1}(G)\)）使得 \(\mathcal{F}(\tau')\geqslant1\) 在 \(\widehat{G}\) 上处处成立。存在 \(\mathcal{M}^{1}(G)\) 的一个非 Hermite 特征标。（利用 a)、b)、练习 3 的 c)，以及 II，第 265 页，练习 15 的 f)。）
\end{enumerate}

\(d)\) 由 \(c\) ) 推出：\(\widehat{\mathrm{G}}\)（它在 \(\mathsf{X}(\mathcal{M}^1 (\mathrm{G}))\) 中是开的）在 \(\mathsf{X}(\mathcal{M}^1 (\mathrm{G}))\) 中不是稠密的。

